\begin{solapartat}
\begin{minipage}[t]{0.6\linewidth}
\punts{0,4} Cal identificar 3 ventres i 4 nodes. Si manquen els dos nodes dels extrems es restarà 0,1~p.

\medskip
\punts{0,3} Del gràfic: $L = 1,5\lambda \Rightarrow \lambda = 0,467\ \mathrm{m}$
\end{minipage}\hfill
\begin{minipage}[t]{0.36\linewidth}
\vspace{0pt}
\figura[1]{sol_fig1.png}
\end{minipage}

\medskip
\destaca{Alternativament}, la longitud d'ona de l'harmònic $n$ és:
\[ \lambda = \frac{2}{n}\,L \]
Per tant, per al tercer harmònic:
\[ \lambda = \frac{2}{3}\,L = 0,467\ \mathrm{m} \]

\punts{0,3} Del gràfic, la distància entre nodes consecutius és $\frac{1}{3}L = 0,233\ \mathrm{m}$

\smallskip
O, també, es pot justificar com que és $\frac{\lambda}{2} = 0,233\ \mathrm{m}$
\end{solapartat}

\begin{solapartat}
\punts{0,4} $f = \dfrac{v}{\lambda} = 660\ \mathrm{Hz}$

\smallskip
\punts{0,3} $T = \dfrac{1}{f} = 1,51\times 10^{-3}\ \mathrm{s}$

\smallskip
\punts{0,3} D'altra banda, el temps que triga un ventre a passar del seu valor màxim al mínim és mig període: $\Delta t = \dfrac{T}{2} = 7,58\times 10^{-4}\ \mathrm{s}$
\end{solapartat}
